% main.tex: "Do Specifications Make LLMs Better Architects? A Pre-Registered Factorial Study of Architectural
% Conformance in LLM-Generated TypeScript"
% Empirical paper draft, target ESEM, MSR or EMSE (about 10-12 pages in this class).
%
% Provenance convention: every number carries a LaTeX comment "% src: <path>" naming the committed file it was read
% from (paths relative to the repository root, branch v1.2e at b7034682). Where a number was derived (a sum, a mean
% over rows, a ratio), the comment says "derived" and how. Nothing here comes from an uncommitted file.
%
% Abbreviations used in the comments:
%   ADR.md      = "Docs/ADR <U+2014> Architectural Decision Records Firewall Tech.md" (its real name contains an em dash,
%                 which this draft does not reproduce)
%   GRID        = results/e1-grid/agg/so5_grid.csv
%   TESTS       = results/e1-grid/agg/so5_tests.csv
%   SIZE        = Docs/DiagnosticRuns/e1-size-analysis-2026-10-10.md
%   STESTS      = results/e1-grid/agg-size/size_tests.csv
%   SFUNC       = results/e1-grid/agg-size/size_functions.csv
%   AGGR        = results/e1-grid/agg-proportional/ahs_by_aggregation.csv
%   JSUM        = results/e1-grid/judge-diagnostics/judge_summary.csv
%   P4          = results/e1-grid/label-plan-p4-api/labelling.json (field "agreement" unless stated)
%   PLAN        = Docs/analysis-plan.md
%   GENP        = Docs/generator-protocol.md
%   TTV         = Docs/threats-to-validity.md
%
% Style rules for editors: no em dashes anywhere (not as a character, not as "---"); every reading is labelled
% pre-registered or post-hoc; do not claim that specifications make architecture worse.

\documentclass[10pt,a4paper]{article}

\usepackage[utf8]{inputenc}
\usepackage[T1]{fontenc}
\usepackage{lmodern}
\usepackage[margin=2.0cm]{geometry}
\usepackage{amsmath}
\usepackage{booktabs}
\usepackage{tabularx}
\usepackage{multirow}
\usepackage{xcolor}
\usepackage{enumitem}
\usepackage{pgfplots}
\pgfplotsset{compat=1.17}
\usetikzlibrary{patterns}
\usepackage[hidelinks]{hyperref}

% \citeneeded{key}: a citation written from general knowledge that must be checked before submission.
\newcommand{\citeneeded}[1]{\cite{#1}\textsuperscript{\textcolor{red}{\tiny[verify]}}}
% \todo{...}: an open item for the authors (also listed in README.md).
\newcommand{\todo}[1]{\textcolor{blue}{[TODO: #1]}}

% Registration-status tags, used on every reading.
\newcommand{\regC}{\textsf{\scriptsize[reg-conf]}}  % pre-registered, confirmatory
\newcommand{\regX}{\textsf{\scriptsize[reg-expl]}}  % pre-registered, exploratory or descriptive
\newcommand{\ph}{\textsf{\scriptsize[post-hoc]}}    % decided after the E1 results were known

\newcommand{\tool}{DaedalusArch}
\newcommand{\ff}[1]{\texttt{\small #1}}
\newcommand{\lev}[1]{\texttt{\small #1}}
\newcommand{\AHSc}{\mathrm{AHS}_{\mathrm{comb}}}
\newcommand{\AHSd}{\mathrm{AHS}_{\mathrm{det}}}
\newcommand{\AHSn}{\mathrm{AHS}_{\mathrm{neur}}}

\setlist{itemsep=1pt,topsep=3pt}

\title{Do Specifications Make LLMs Better Architects?\\[2pt]
\large A Pre-Registered Factorial Study of Architectural Conformance in LLM-Generated TypeScript}
\author{\todo{author list and affiliations}}
\date{Draft, 2026-10-11}

\begin{document}
\maketitle

% =====================================================================================================================
\begin{abstract}
\textbf{Context.} Spec-driven development asks coding agents to follow a written architecture, but it is not known
whether a specification changes the architecture an agent builds, or by how much a machine-readable specification
differs from prose. \textbf{Objective.} We test whether the specification level (none, a short prose description,
or a full Architecture-as-Code (AaC) specification) and the model change the architectural conformance of
TypeScript services generated by LLM agents. \textbf{Method.} A pre-registered factorial experiment: three Claude
models $\times$ three specification levels $\times$ two tasks $\times$ three runs (54 cells), plus an exploratory
18-cell arm with one OpenAI model through the Codex CLI. Each project is scored by a neuro-symbolic conformance
instrument against a clean-architecture specification; an LLM judge contributes 8\% of the score and is checked
against a model panel from another vendor. \textbf{Results.} Specification level has a registered effect on the
Architectural Health Score (Holm $p = 0.0003$), but in the reverse direction: projects built without a specification
score best. A post-hoc analysis explains why. Without a specification, models build small, flat projects (about 9
files against 27 to 30) where coupling and dependency-inversion rules have little to check; findings per file do not
differ between levels. Two rules, module fan-out and dependency inversion, carry the gap. Between the two
specification styles, which produce projects of the same size, the full AaC specification beats the short prose one
(Cliff's $\delta = 0.52$; $0.47$ after a size adjustment), so the vague specification is the worst case. Model size
has no effect. The exploratory Codex arm prefers the prose specification. The judge changes no verdict at its 8\%
weight; a registered self-preference check is positive ($p = 0.014$) but is carried by the weakest model's cells;
on a 36-item sample the judge's weighted agreement with the panel is 0.982 (AC1 0.98; 34 of 36 items agree), and both disagreements are
lenient; replacing the judge model on 100 recorded units changes no verdict ($\kappa = 0.77$).
% src: P4 headline row (percentAgreement 0.982301, AC1 0.981985); results/e1-grid/label-plan-p4-api/labels.json joined with results/e1-grid/label-plan-p4/judge-verdicts.json (36 decided E1 items: 34 agree, 2 judge pass / panel fail; derived)
% src: results/judge-robustness/summary.json (units 100, kappa 0.766355, registeredVerdictChanges 0); ADR.md ADR-038
\textbf{Conclusion.} A specification changes how much structure an agent builds, and scores that reward
small structures can make that change look like a loss. At equal size, a precise AaC specification gives more
conformant code than prose.
\end{abstract}

% =====================================================================================================================
\section{Introduction}\label{sec:intro}

Coding agents built on large language models (LLMs) now generate whole services from a task description
\citeneeded{chen2021codex,jimenez2024swebench}. A growing practice, spec-driven development, puts a written
specification in front of the agent so that the generated code follows an intended design
\citeneeded{speckit2025,kiro2025}. For architecture this raises a measurable question: does a specification change
the architecture an agent builds, and does the form of the specification matter? A specification can be a short
prose paragraph (``use clean architecture; dependencies point inward'') or a machine-readable Architecture-as-Code
(AaC) document that lists layers, roles, fitness functions and thresholds \citeneeded{ford2017evolutionary}. The two
cost the author very different amounts of effort, and practitioners need to know whether the extra precision pays.

Functional correctness of generated code is well studied \citeneeded{chen2021codex,liu2023evalplus}; architectural
conformance is not. Architecture conformance checking has a long history for human-written code
\citeneeded{murphy1995reflexion,knodel2007comparison,passos2010static}, and architecture erosion is a known cost
\citeneeded{desilva2012erosion}. Generated code adds two new questions: whether the specification the agent reads
changes conformance, and whether an LLM used as a judge of that code is biased toward code from its own model family
\citeneeded{panickssery2024self,wataoka2024selfpref}.

We report E1, a pre-registered factorial experiment. Three Claude models generate two Express services at three
specification levels, three times each; an added arm runs one OpenAI model through the Codex CLI. A neuro-symbolic
conformance instrument, \tool{}, described in a companion paper \cite{daedalus2026tool}, scores every project
against a clean-architecture AaC specification. The analysis plan, the seeds, the generator protocol and the
prompt templates were registered and hash-checked before any cell ran. We ask:

\begin{description}[leftmargin=1.2cm,style=sameline]
  \item[RQ1] Does the specification level (none, minimal prose, full AaC) change architectural conformance?
  \item[RQ2] Do model capability (Haiku, Sonnet, Opus) and vendor (the exploratory Codex arm) matter?
  \item[RQ3] Which failure families drive the differences?
  \item[RQ4] What does an LLM judge add to the conformance score, and is it biased?
\end{description}

\paragraph{Findings in brief.} The registered test finds a specification-level effect, but the projects built with
no specification score best. Our post-hoc size analysis shows that this is an effect of the specification on how
much structure the agent builds: without a specification the agents write about 9 files, with one they write 27 to
30, and two rules that need structure to fire (module fan-out and dependency inversion) almost never fire on the
small projects. Findings per file are the same at every level. Between the two specification styles, which give
projects of the same size, the full AaC specification beats the prose one. The model has no effect. The LLM judge
does not move any verdict, a registered self-preference check is positive but confounded, and a cross-vendor panel
agrees with the judge on almost every sampled item.

\paragraph{Contributions.}
\begin{enumerate}
  \item A pre-registered $3 \times 3 \times 2 \times 3$ experiment on specification level and model for agentic
        generation of TypeScript services, with an exploratory second-vendor arm (Sections~\ref{sec:design}
        and~\ref{sec:analysis}).
  \item Evidence that a registered conformance score can reverse the apparent effect of a specification through
        project size, and a decomposition that locates the reversal in two rules (Section~\ref{sec:results}).
  \item A size-matched comparison of prose and AaC specifications that favours AaC, and practical guidance on what to
        put in a specification (Section~\ref{sec:discussion}).
  \item A registered analysis of an LLM judge inside a conformance score: its effect on verdicts, a self-preference
        check, and judge-versus-panel agreement with a judge from another vendor.
  \item A replication package with every generated tree's hash, every judge and labeller call as a replayable
        cassette, and scripts that regenerate every table and figure (Section~\ref{sec:conclusion}).
\end{enumerate}

Throughout, every reading carries one of three tags: \regC{} for the registered confirmatory tests, \regX{} for
registered exploratory or descriptive outcomes, and \ph{} for analyses decided after the E1 results were known.

% =====================================================================================================================
\section{Background}\label{sec:background}

\paragraph{Architecture as code and fitness functions.} An architectural fitness function is an executable check of
an architectural property, such as dependency direction or coupling \citeneeded{ford2017evolutionary,ford2022evolutionary2}.
Architecture-as-Code expresses the intended architecture as a versioned, machine-readable document from which such
checks are derived. Clean architecture \citeneeded{martin2017clean} is a common target: a domain layer with
entities and repository interfaces, an application layer with use cases, and an infrastructure layer with
controllers and adapters, with dependencies pointing inward.

\paragraph{Conformance checking.} Static conformance checkers compare the implemented dependency structure with an
intended one \citeneeded{murphy1995reflexion,koschke2003hierarchical,terra2009dcl}. In TypeScript, tools such as
dependency-cruiser express path-based import rules \citeneeded{depcruise}. \tool{} \cite{daedalus2026tool} compiles an
AaC specification into graph queries over a property graph of the project and adds rules that path patterns cannot
express (dependency inversion, coupling metrics, SOLID proxies, conventions), plus an LLM judge for two semantic
properties. Section~\ref{sec:instrument} summarises what E1 needs from it.

\paragraph{Spec-driven generation.} Spec-driven tools ask the developer to write a specification, then have an
agent plan and implement against it \citeneeded{speckit2025,kiro2025}. Whether the form of the specification matters
has, to our knowledge, not been tested under a registered design \citeneeded{dhar2024adr,jahic2024practice}.

\paragraph{LLM judges.} LLMs are widely used as evaluators. They show position and verbosity biases
\citeneeded{wang2023fair,zheng2023judging} and can prefer their own generations
\citeneeded{panickssery2024self,wataoka2024selfpref}. In E1 the judge is a Claude model and three of the four
generators are Claude models, so self-preference is a registered concern.

% =====================================================================================================================
\section{Study Design}\label{sec:design}

\subsection{Factors and grid}

The registered grid crosses three factors and replicates each cell three times
(Table~\ref{tab:grid}). % src: experiments/e1-grid/generator-plan.json; experiments/e1-grid/plan.json
\emph{Model} has three Claude levels: \texttt{claude-haiku-4-5}, \texttt{claude-sonnet-5-5} and
\texttt{claude-opus-5-5}, run through the Claude Code CLI (version 2.1.294).
% src: experiments/e1-grid/generator-plan.json (adapters); GENP §10 (CLI 2.1.294)
\emph{Specification level} has three levels (Section~\ref{sec:levels}). \emph{Task} is a blocking factor with two
levels, \texttt{task-management} and \texttt{order-fulfilment}. The three runs are replicates: the Claude Code
headless mode exposes no temperature or seed, so run-to-run variation is part of the within-cell variance
(a declared deviation). % src: PLAN §6 ("The generator ... has no temperature or seed setting"); TTV TV-78
This gives $3 \times 3 \times 2 \times 3 = 54$ Claude cells. % src: GENP §9 ("E1 = 3 models x 2 tasks x 3 levels x 3 runs = 54")

\begin{table}[t]
\centering\small
\caption{The E1 grid. All four arms share the prompt templates, the skeleton, the type-check of record and the
evaluation. The Codex arm was added before any of its cells ran and is exploratory.}
\label{tab:grid}
\begin{tabularx}{\linewidth}{@{}lXr@{}}
\toprule
Factor & Levels & $n$ \\
\midrule
Model (registered) & \texttt{claude-haiku-4-5}, \texttt{claude-sonnet-5-5}, \texttt{claude-opus-5-5} (Claude Code CLI 2.1.294) & 3 \\
Model (exploratory arm) & \texttt{gpt-5.6-terra} (Codex CLI 0.162.1, reasoning effort \texttt{medium}) & 1 \\
Specification level & \lev{none}, \lev{minimal-prose}, \lev{full-aac} & 3 \\
Task (block) & \texttt{task-management}, \texttt{order-fulfilment} & 2 \\
Runs (replicates) & 0, 1, 2 & 3 \\
\midrule
Cells & 54 Claude + 18 Codex & 72 \\
\bottomrule
\end{tabularx}
% src: experiments/e1-grid/generator-plan.json; experiments/e1-grid/generator-plan-codex.json (cliVersion 0.162.1,
%      reasoningEffort medium); experiments/e1-grid/plan.json (e1.models, specLevels, tasks, runs 3); GENP §10
\end{table}

\subsection{Tasks and specification levels}\label{sec:levels}

Both tasks ask for an Express REST API with in-memory storage. \texttt{task-management} has projects, tasks with a
status, assignment and filtering, and two business rules (a task cannot be completed twice; only existing users can
be assignees). \texttt{order-fulfilment} has orders that reserve stock, a five-state life cycle in a fixed order, and
cancellation that releases stock. % src: scripts/generator/prompts/minimal-prose.md (task sections)
Every prompt has the same four parts: the task brief, one sentence that fixes the folder layout (\texttt{src/domain},
\texttt{src/application}, \texttt{src/infrastructure}), an environment paragraph (available packages, no network,
the compiler options and the exact type-check command), and a level block. % src: GENP §5
The level block differs:
\begin{itemize}
  \item \lev{none}: no block. The agent still sees the three folder names.
  \item \lev{minimal-prose}: about 140 words of prose derived from the clean-architecture preset, naming every layer
        and role and the dependency direction, with no rule identifier and no number. % src: GENP §5
  \item \lev{full-aac}: the clean-architecture preset verbatim, a 307-line YAML document with layers, roles, fitness
        functions, thresholds and weights. % src: GENP §5 ("the preset verbatim"); presets/clean-architecture.yaml (wc -l = 307)
\end{itemize}
The \lev{full-aac} preset carries the same rules, thresholds and weights as the specification the instrument
evaluates with (\texttt{specs/clean-arch.yaml}); apart from comments and quoting, the two files differ only in a
\texttt{default\_exclude\_paths} block (\texttt{node\_modules}, \texttt{dist}, spec and test files) that the preset carries.
% src: derived, diff of presets/clean-architecture.yaml and specs/clean-arch.yaml with comment and blank lines removed
% (besides comments and quoting, only the preset's default_exclude_paths block, 5 lines, differs; spec_version 1.0.0 in both; 27 FF ids, weights and thresholds equal); experiments/e1-grid/plan.json (specPath specs/clean-arch.yaml)
At \lev{full-aac} the agent therefore sees the test it will be scored against. This confound was declared before
the run (Section~\ref{sec:threats}). % src: PLAN §10 B5
The three template files differ in length by more than an order of magnitude (21, 25 and 637 lines).
% src: derived, wc -l scripts/generator/prompts/{none,minimal-prose,full-aac}.md (21, 25, 637; two task sections each)

\subsection{Generation protocol}

The generator protocol was frozen on 2026-10-08, after confinement probes and a three-cell pilot that is not part of
E1. % src: GENP header ("FROZEN 2026-10-08") and §10
Each cell is one headless agent call in a fresh directory outside the repository, with the tools Read, Write, Edit,
Glob, Grep and Bash, where Bash is restricted by an allow rule to one exact type-check command.
% src: GENP §1
The directory holds a read-only skeleton with exactly four pinned packages (\texttt{express} 5.2.1,
\texttt{@types/express} 5.0.6, \texttt{@types/node} 22.19.15, \texttt{typescript} 5.9.3).
% src: GENP §4 (skeleton table)
The type-check of record uses a pinned strict tsconfig kept outside the tree. % src: GENP §3
The protocol assigns each run one status in a fixed order; a run that fails the type-check is
\texttt{failed-typecheck}, is recorded and is never replaced. % src: GENP §7 (status order; "Failed generations are recorded and never replaced")
A run must also be model-valid: the pinned model must produce the largest share of output tokens.
% src: GENP §6 (model-usage rule)
The 54 Claude cells ran in an order drawn from a registered seed (20261008), blocked by run index.
% src: experiments/e1-grid/generator-plan.json (orderSeed 20261008); GENP §9
Usage-limit pauses consume no attempt and resume the same cell. % src: GENP §8

\paragraph{The Codex arm.} After the Claude grid was registered and before any E1 cell ran, an 18-cell arm was added:
\texttt{gpt-5.6-terra} through the Codex CLI 0.162.1, the strongest coding model generally available on the
author's plan, with its own plan file and order seed (20261010). % src: ADR.md ADR-029 items 1-3; experiments/e1-grid/generator-plan-codex.json
It shares the prompt templates, skeleton, tsconfig, status order and evaluation, but its harness, tools and
confinement mechanism differ: the Codex agent may run any command inside an OS sandbox, while the Claude agent may run
only the type-check. % src: PLAN §10 B10; ADR.md ADR-029 item 4
Every registered test stays on the 54 Claude cells; the Codex cells enter only the vendor families.
% src: PLAN §6.1 ("Registered analyses unchanged")

\subsection{Pre-registration}

The analysis plan, generator plans, prompt templates and seeds are registered artefacts whose hashes are recorded
in \texttt{corpus/prereg.json}; a run is refused when a registered file has changed. % src: PLAN header; GENP §11
Every amendment is a dated version bump with a reason. Table~\ref{tab:timeline} lists those that touch E1.

\begin{table}[t]
\centering\small
\caption{Registration timeline for E1 (all dates 2026). ``Before'' and ``after'' refer to the first E1 generation.}
\label{tab:timeline}
\begin{tabularx}{\linewidth}{@{}llX@{}}
\toprule
Date & Version & Change \\
\midrule
10-08 & v1, P-2 & Analysis plan; co-primary $\AHSd$ for the model effect; directional self-preference check; descriptive pairwise intervals (before) \\
10-09 & v4 (P-U6) & Stratified permutation; size and density columns; open-coding procedure (before) \\
10-10 & v14--v17 & Proportional judge aggregation, judge diagnostics, judge-weighted variant as registered sensitivity analyses (before) \\
10-10 & v18 & Codex arm, vendor contrast and vendor self-preference check (before any Codex cell) \\
10-10 & v20 & ADR-032: pinned tsconfig for extraction, operational and post-hoc (after generation, during the judge pass) \\
10-11 & v22, v23 & ADR-034 and ADR-035: labeller route for the judge-vs-panel sample (after E1 scoring, before labelling) \\
10-11 & v25 & ADR-038: judge model robustness, Sonnet against Opus on recorded judge units (after E1 scoring, before any Sonnet call) \\
\bottomrule
\end{tabularx}
% src: PLAN header amendments (v1, P-2, P-U6 v4, v14, v15, v17, v18); ADR.md ADR-032 (prereg v20), ADR-034, ADR-035 (prereg v23);
%      aidlc-docs/audit.md entries "go E1", "P4 labelling stop and ADR-034", "P4 stop rule fired; ADR-035"
%      ADR.md ADR-038 (prereg v25, "registered before it runs", date 2026-10-11)
\end{table}

% =====================================================================================================================
\section{Instrument}\label{sec:instrument}

\tool{} is described and validated in the companion paper \cite{daedalus2026tool}; we summarise what E1 uses.

\paragraph{Pipeline.} The instrument extracts an Architectural Property Graph (APG) from the TypeScript project
(files, classes, interfaces, imports, re-exports, constructor injection, inheritance, calls), loads it into a graph
database and runs one compiled query per symbolic fitness function. Two neural functions, \ff{FF-N01}
(architectural integrity) and \ff{FF-N02} (intent alignment), are judged by an isolated LLM on a capped sample of
units; each unit is judged three times and decided by majority. % src: papers/paper2-tool/main.tex §3.4 ("The isolated judge")
In E1 the judge is \texttt{claude-opus-5-5}. % src: experiments/e1-grid/plan.json (judge.model)
The clean-architecture specification declares 27 fitness functions. % src: derived, grep -c "id: FF-" specs/clean-arch.yaml = 27
E1 ran in full mode (symbolic and neural) with instrument version 2. % src: SIZE header ("Instrument: v2 (B8)"); ADR.md ADR-029 / director report §4.8.1

\paragraph{Score.} Each function belongs to one of seven dimensions. A dimension's Architectural Violation Ratio
$\mathrm{AVR}_d$ is the confidence-weighted share of its executed functions that fail, and
\begin{equation}
  \mathrm{AHS} = \sum_d w_d\,(1 - \mathrm{AVR}_d).
\end{equation}
$\AHSd$ uses the five symbolic dimensions; $\AHSc$ adds the two judge dimensions. The full-mode weights are
structural 0.32, coupling 0.18, pattern 0.27, SOLID 0.10, convention 0.05, semantic 0.04 and integrity 0.04, so the
judge carries $W_n = 0.08$ in total. The verdict is \emph{pass} at $\geq 0.80$, \emph{warning} at $\geq 0.65$,
\emph{soft-block} at $\geq 0.50$ and \emph{hard-block} below.
% src: specs/clean-arch.yaml (scoring.full_mode_weights, scoring.thresholds); papers/paper2-tool/main.tex §3.5
A neural function fails under the registered rule when strictly more than half of its valid units fail.
% src: PLAN §12.1
Two properties of the score matter for what follows. First, a function counts once in its dimension when it fires,
however many findings it produces: the AHS measures \emph{which} functions fire, not how many findings there are.
% src: SIZE §5 ("Why counts do not move the AHS")
Second, with $W_n = 0.08$ the judge can lower $\AHSc$ by at most 0.08 relative to its judge-blind value and can
never raise it. % src: papers/paper2-tool/main.tex §3.5; ADR.md ADR-028 review fixes

\paragraph{Validity evidence from the companion study.} On 73 seeded held-out violations in seven open-source
projects, the symbolic path reaches recall 0.918 and labelled precision 0.957. % src: ADR.md ADR-033 item 1
Precision on unseeded real findings is much less certain: 0.328 under the registered rule (\texttt{uncertain} labels
counted as not real; Wilson-Kish 95\% interval 0.132 to 0.609), 0.608 with \texttt{uncertain} excluded and 0.789
with \texttt{uncertain} counted as real, the last two post-hoc. % src: results/so4-heldout/agg/precision_baseline.csv (v2 overall); ADR.md ADR-033 item 4
We return to this range in Section~\ref{sec:threats}.

% =====================================================================================================================
\section{Analysis}\label{sec:analysis}

\paragraph{Unit and outcomes.} The unit of analysis is one valid cell: an accepted evaluation of a generation with
status \texttt{ok}. The primary outcome is $\AHSc$ \regC{}. For the model effect, $\AHSd$ is co-primary \regC{},
because a model effect seen only in the judge-bearing field could be judge self-preference.
% src: PLAN §6 (Unit of analysis; Primary outcome; Co-primary)
Registered exploratory outcomes \regX{} include $\AHSn$, each dimension's AVR, the judge fail share, the valid-generation
yield and deterministic violations per 1\,000 lines of code (KLOC). Lines of code, all violations per KLOC and the
latencies are registered descriptive columns. % src: PLAN §3 (SO5 row) and "SO5 size, density and latency definitions"

\paragraph{Tests.} Each outcome gets three permutation tests (10\,000 permutations, registered seed 5103): model,
specification level and their interaction. The statistic is the between-level sum of squares of the tested factor;
model labels are permuted within (task, specification level) and specification labels within (task, model), so an
unequal number of valid runs per cell cannot carry one factor's effect into the other's test. The three tests of an
outcome form one Holm family \citeneeded{holm1979}, at $\alpha = 0.05$.
% src: PLAN §6 (Tests; Strata; Holm families; Significance level); experiments/e1-grid/plan.json (seeds.permutation 5103)
Pairwise differences carry a cluster-bootstrap 95\% interval (runs resampled within cells, seed 5102) and Cliff's
$\delta$ \citeneeded{cliff1993}. The intervals have no multiplicity control and are descriptive only; inference
rests on the Holm-corrected $p$ values and $\delta$. % src: PLAN §6 (Effect sizes); experiments/e1-grid/plan.json (seeds.bootstrap 5102)
We use permutation tests because 47 cells with unequal replication give no reason to trust a parametric model
\citeneeded{good2005permutation}.

\paragraph{Self-preference checks.} The registered directional check \regC{} computes, per valid cell,
$d = \AHSn - \AHSd$ and tests one-sided whether mean $d$ is larger for cells written by the judge's own model
(\texttt{claude-opus-5-5}) than for the other Claude cells (family of one, $\alpha = 0.05$). The registered vendor
check \regC{} does the same for Claude-written against GPT-written cells. A positive result is to be reported as
evidence of self-preference; a null vendor result is not evidence of its absence, because power is low.
% src: PLAN §6 (Directional self-preference check) and §6.1 (Vendor self-preference check)

\paragraph{The judge readings.} Three judge sensitivity analyses were registered before E1 \regX{}: a proportional
aggregation (a neural dimension's AVR equals the weighted share of failing units), an any-fail companion, and a
judge-weighted variant that gives the judge dimensions 2/7 of the weight. They are descriptive, with no test.
% src: PLAN §12.2, §12.3, §12.5; AGGR column reading = pre-registered on all 65 rows

\paragraph{Judge versus panel.} A stratified sample of judge units (population P4) is labelled by a panel model of
another vendor, \texttt{gemini-3.1-pro-preview} at temperature 0 with tools off, two runs per item, with the same
persona and prompts as the other labelling populations. % src: ADR.md ADR-035 items 1-2; Docs/labeller-route.md
The registered headline is weighted judge-vs-panel agreement on the E1 units, with Cohen's $\kappa$
\citeneeded{cohen1960} and Gwet's AC1 \citeneeded{gwet2008} and weighted item-bootstrap intervals, reported per
source and per generator model. % src: PLAN §3 (SO3 neural row); P4 (ciMethod weighted-item-bootstrap)
The route was changed once before labelling the reported answers: the first route (a CLI agent) stopped on a
registered rule because it used a web-search tool on 2 of 10 items, and the author moved P4 to the API with tools
off (ADR-034, ADR-035). The 8 answers of the stopped route are kept apart and are not used here.
% src: ADR.md ADR-035 Context and item 4; aidlc-docs/audit.md "P4 stop rule fired"

\paragraph{The post-hoc size analysis.} After the E1 results were known, we ran a diagnostic analysis on the
committed data, with no new run or model call. % src: SIZE header and §1
It first reproduces all 27 registered rows of the primary, co-primary and density families to the last printed
digit. % src: SIZE §2 ("All 27 registered rows ... equal the committed so5_tests.csv rows"); STESTS column reproduces_registered
It then adds tests on file count, LOC, findings per file, density without one rule, and $\AHSc$ residualised on
$\ln(\mathrm{LOC})$, each with the registered procedure in its own Holm family of three, with no correction across
families; and a split of the AHS gap by function. Every one of these readings is \ph{}.
% src: SIZE §3 ("The post-hoc rows each use their own Holm family of three") and §8

% =====================================================================================================================
\section{Results}\label{sec:results}

\subsection{Yield and the shape of the data}

Of the 72 cells, 65 were evaluated. The 7 others are \texttt{claude-haiku-4-5} generations that failed the
type-check of record (3 at \lev{none}, 2 at \lev{minimal-prose}, 2 at \lev{full-aac}); they are recorded and never
replaced. % src: GRID (status, generation_status; derived counts by spec_level)
Opus, Sonnet and the Codex model produced 18 of 18 valid cells each, giving 47 valid Claude cells and 18 valid Codex
cells. % src: GRID (derived counts of status = accepted by requested_model_id)
Valid-generation yield differs by model \regX{} (Holm $p = 0.004$; Haiku minus Opus $-0.39$) but not by
specification level (Holm $p = 1.000$). % src: TESTS family secondary:valid_generation_yield (model p_holm 0.004200; spec-level p_holm 1.000000; haiku-opus -0.388889)

\textbf{Every evaluated cell passes.} The lowest $\AHSc$ is 0.852 and the highest 0.988, against the 0.80 pass
threshold, and the structural dimension (dependency direction and cycles) is clean in every cell.
% src: GRID (derived min/max of ahs_combined over accepted rows); TESTS family secondary:avr_structural (all statistics 0)
E1 is therefore a study of differences inside the pass band. No agent at any level produced an inward-dependency
violation that the instrument detects.

\subsection{RQ1: Specification level}\label{sec:rq1}

\paragraph{Registered result.} On the 47 valid Claude cells, specification level has an effect on $\AHSc$ \regC{}
(Holm $p = 0.0003$), and so does the model $\times$ level interaction (Holm $p = 0.026$).
% src: TESTS family primary:ahsCombined (spec-level p_holm 0.000300; model×spec-level p_holm 0.026397)
The direction is the reverse of the hypothesis: \lev{none} scores best (mean 0.951), \lev{minimal-prose} worst
(0.887) and \lev{full-aac} in between (0.920) (Table~\ref{tab:levels}, Figure~\ref{fig:ahs}).
% src: SIZE §3 (per-level table, "all" rows); derived check: mean ahs_combined of accepted Claude rows in GRID by spec_level
All three pairwise contrasts are large (Table~\ref{tab:pairs}). The same order holds for $\AHSd$ (\regX{} for the
level effect, Holm $p = 0.0003$). % src: TESTS family co-primary:ahsDeterministic (spec-level p_holm 0.000300, exploratory true)

\begin{table}[t]
\centering\small
\caption{Per-level means over valid cells. Density is deterministic violations per KLOC. Claude rows pool the three
models; Codex rows are the exploratory arm. Tags give the status of each column.}
\label{tab:levels}
\begin{tabular}{@{}llrrrrrr@{}}
\toprule
Arm & Level & Valid & $\AHSc$ & Files & LOC & Viol./KLOC & Viol./file \\
 & & & \regC{} & desc. & \regX{} desc. & \regX{} & \ph{} \\
\midrule
Claude & \lev{none}          & 15 & 0.951 &  9.1 & 411 & 27.0 & 1.14 \\
Claude & \lev{minimal-prose} & 16 & 0.887 & 27.4 & 551 & 58.1 & 1.15 \\
Claude & \lev{full-aac}      & 16 & 0.920 & 29.6 & 600 & 57.0 & 1.15 \\
\midrule
Codex  & \lev{none}          &  6 & 0.926 &  6.3 & 254 & 32.0 & 1.31 \\
Codex  & \lev{minimal-prose} &  6 & 0.957 & 10.5 & 274 & 43.5 & 1.13 \\
Codex  & \lev{full-aac}      &  6 & 0.942 & 13.3 & 320 & 48.4 & 1.19 \\
\bottomrule
\end{tabular}
% src: SIZE §3 (Claude "all" rows) and §7 (Codex table); results/e1-grid/agg-size/size_levels.csv;
%      derived check of Valid, AHS, Files, LOC from GRID (accepted rows, means by arm and spec_level)
\end{table}

\begin{table}[t]
\centering\small
\caption{Specification-level contrasts on the Claude arm (row level minus column level). Differences are means with
the descriptive cluster-bootstrap 95\% interval; $\delta$ is Cliff's delta. The LOC-adjusted row is \ph{}.}
\label{tab:pairs}
\setlength{\tabcolsep}{4pt}
\begin{tabular}{@{}lccc@{}}
\toprule
Outcome & full-aac $-$ none & minimal-prose $-$ none & full-aac $-$ minimal-prose \\
\midrule
$\AHSc$ \regC{}            & $-0.031$, $\delta\,{-0.54}$ & $-0.064$, $\delta\,{-0.93}$ & $+0.033$, $\delta$ 0.52 \\
\quad 95\% interval (desc.) & [$-0.044$, $-0.019$] & [$-0.071$, $-0.058$] & [0.019, 0.047] \\
$\AHSc$ resid.\ $\ln$ LOC \ph{} & $-0.023$, $\delta\,{-0.28}$ & $-0.057$, $\delta\,{-0.87}$ & $+0.034$, $\delta$ 0.47 \\
\quad 95\% interval (desc.) & [$-0.036$, $-0.010$] & [$-0.064$, $-0.051$] & [0.020, 0.049] \\
Files \ph{}                & $+20.5$, $\delta$ 1.00 & $+18.3$, $\delta$ 1.00 & $+2.2$, $\delta$ 0.12 \\
Viol./KLOC \regX{}         & $+30.0$, $\delta$ 0.97 & $+31.1$, $\delta$ 0.97 & $-1.1$, $\delta\,{-0.09}$ \\
Viol./file \ph{}           & $+0.010$, $\delta$ 0.08 & $+0.013$, $\delta$ 0.00 & $-0.003$, $\delta\,{-0.02}$ \\
\bottomrule
\end{tabular}
% src: TESTS family primary:ahsCombined (spec-level pairs: -0.031333 [-0.043646,-0.018958] d -0.541667;
%      -0.064083 [-0.070833,-0.057583] d -0.925; +0.032750 [0.019123,0.046875] d 0.519531);
%      STESTS outcomes ahsCombined_resid_lnloc, files, violations_deterministic_per_file; TESTS family secondary:violations_per_kloc;
%      SIZE §3 test table
\end{table}

\begin{figure}[t]
\centering
% Generated by papers/paper1-e1/figures/make_figures.py from results/e1-grid/agg/so5_grid.csv. Do not edit by hand.
\definecolor{oiBlue}{HTML}{0072B2}
\definecolor{oiOrange}{HTML}{E69F00}
\definecolor{oiGreen}{HTML}{009E73}
\definecolor{oiPink}{HTML}{CC79A7}
\begin{tikzpicture}
\begin{axis}[width=\linewidth, height=4.9cm, ymin=0.84, ymax=1.0, xmin=-0.5, xmax=2.5,
  xtick={0,1,2}, xticklabels={none, minimal-prose, full-aac}, ylabel={$\mathrm{AHS}_{\mathrm{comb}}$},
  ytick={0.85,0.90,0.95,1.00}, yticklabel style={/pgf/number format/fixed,/pgf/number format/precision=2},
  grid=major, grid style={gray!20}, axis line style={gray!60}, tick style={gray!60},
  legend style={font=\scriptsize, at={(0.5,1.02)}, anchor=south, legend columns=4, draw=none},
  every axis plot/.append style={only marks}]
\addplot[mark=*, mark size=1.3pt, draw=oiBlue, fill=oiBlue, fill opacity=0.35, draw opacity=0.6, forget plot]
  coordinates {(-0.305,0.914) (-0.270,0.884) (-0.235,0.988) (0.695,0.900) (0.730,0.912) (0.765,0.928) (0.695,0.900) (1.695,0.882) (1.730,0.929) (1.765,0.900) (1.695,0.869)};
\addplot[mark=*, mark size=3pt, draw=black, fill=oiBlue, line width=0.4pt]
  coordinates {(-0.270,0.9287) (0.730,0.9100) (1.730,0.8950)};
\addlegendentry{Haiku 4.5}
\addplot[mark=square*, mark size=1.3pt, draw=oiOrange, fill=oiOrange, fill opacity=0.35, draw opacity=0.6, forget plot]
  coordinates {(-0.125,0.957) (-0.090,0.957) (-0.055,0.957) (-0.125,0.957) (-0.090,0.957) (-0.055,0.957) (0.875,0.852) (0.910,0.852) (0.945,0.852) (0.875,0.928) (0.910,0.928) (0.945,0.882) (1.875,0.957) (1.910,0.912) (1.945,0.882) (1.875,0.957) (1.910,0.882) (1.945,0.928)};
\addplot[mark=square*, mark size=3pt, draw=black, fill=oiOrange, line width=0.4pt]
  coordinates {(-0.090,0.9570) (0.910,0.8823) (1.910,0.9197)};
\addlegendentry{Sonnet 5.5}
\addplot[mark=triangle*, mark size=1.3pt, draw=oiGreen, fill=oiGreen, fill opacity=0.35, draw opacity=0.6, forget plot]
  coordinates {(0.055,0.957) (0.090,0.957) (0.125,0.957) (0.055,0.957) (0.090,0.957) (0.125,0.957) (1.055,0.882) (1.090,0.882) (1.125,0.852) (1.055,0.882) (1.090,0.882) (1.125,0.882) (2.055,0.882) (2.090,0.957) (2.125,0.912) (2.055,0.957) (2.090,0.957) (2.125,0.957)};
\addplot[mark=triangle*, mark size=3pt, draw=black, fill=oiGreen, line width=0.4pt]
  coordinates {(0.090,0.9570) (1.090,0.8770) (2.090,0.9370)};
\addlegendentry{Opus 5.5}
\addplot[mark=diamond*, mark size=1.3pt, draw=oiPink, fill=oiPink, fill opacity=0.35, draw opacity=0.6, forget plot]
  coordinates {(0.235,0.896) (0.270,0.929) (0.305,0.924) (0.235,0.924) (0.270,0.924) (0.305,0.957) (1.235,0.957) (1.270,0.957) (1.305,0.957) (1.235,0.957) (1.270,0.957) (1.305,0.957) (2.235,0.957) (2.270,0.924) (2.305,0.957) (2.235,0.957) (2.270,0.957) (2.305,0.900)};
\addplot[mark=diamond*, mark size=3pt, draw=black, fill=oiPink, line width=0.4pt]
  coordinates {(0.270,0.9257) (1.270,0.9570) (2.270,0.9420)};
\addlegendentry{GPT-5.6-terra (Codex, expl.)}
\draw[gray!70, dashed] (axis cs:-0.5,0.80) -- (axis cs:2.5,0.80);
\end{axis}
\end{tikzpicture}

\caption{$\AHSc$ per valid cell (small marks) and per model and level (large marks), by specification level. The
dashed line is the pass threshold, below the plotted range. Text alternative: Table~\ref{tab:levels} gives the
per-level means; on Claude models \lev{none} is highest and \lev{minimal-prose} lowest; on the Codex model
\lev{minimal-prose} is highest and \lev{none} lowest. Generated from \texttt{results/e1-grid/agg/so5\_grid.csv} by
\texttt{figures/make\_figures.py}.}
\label{fig:ahs}
\end{figure}

\paragraph{Without a specification, agents build small, flat projects \ph{}.} On the Claude arm the mean file count
is 9.1 at \lev{none}, 27.4 at \lev{minimal-prose} and 29.6 at \lev{full-aac}; \lev{none} differs from each other
level with $\delta \geq 0.99$ (Holm $p = 0.0003$). % src: SIZE §3; STESTS outcome files (spec-level p_holm 0.000300; delta 0.995833, 1.000000)
The \lev{none} files are about twice as long (about 45 against 20 LOC per file). % src: SIZE §6 ("45 vs 20 LOC per file")
The agents are told the three folder names at every level, so the \lev{none} projects are layered, but thinly: few
files, few interfaces and few modules per layer.

\paragraph{Findings per file do not differ \ph{}.} Deterministic findings per file are 1.14, 1.15 and 1.15 across
the levels (Holm $p = 0.95$), and still show no level effect without the test-file-pairing rule (Holm $p = 0.77$).
% src: SIZE §3 (Viol./file column); STESTS outcomes violations_deterministic_per_file (spec-level p_holm 0.948705) and
%      violations_per_file_excl_fpat_pattern (p_holm 0.774823)
Per KLOC, the registered exploratory density, \lev{none} has about half the findings of the other levels (27.0
against 58.1 and 57.0; Holm $p = 0.0003$) \regX{}, because its files are longer.
% src: TESTS family secondary:violations_per_kloc (spec-level p_holm 0.000300); SIZE §3
The density result follows the AHS result in direction, but per file the three levels look the same.

\paragraph{Size is on the causal path \ph{}.} The specification causes the size difference, so size is a mediator,
not an outside confounder, and no adjustment for file count is identifiable. % src: SIZE §6 ("Size as a mechanism, not a confounder")
Adjusting the AHS for $\ln(\mathrm{LOC})$ leaves the level effect (Holm $p = 0.0003$) and shrinks the
\lev{full-aac} versus \lev{none} gap from $\delta = -0.54$ to $-0.28$ (Table~\ref{tab:pairs}).
% src: STESTS outcome ahsCombined_resid_lnloc (spec-level p_holm 0.000300); SIZE §3
Within a level, size does not predict the AHS consistently: at \lev{full-aac} the larger projects score
\emph{higher} (files-AHS Pearson $r = +0.59$), while at the other levels the correlation is weakly negative
($r = -0.26$ and $-0.25$; Figure~\ref{fig:files}). These correlations are descriptive.
% src: SIZE §4 (table); results/e1-grid/agg-size/size_correlations.csv


\paragraph{The size-matched contrast: AaC beats prose \regC{} / \ph{}.} \lev{full-aac} and \lev{minimal-prose} give
projects of nearly the same size (files $\delta = 0.12$, LOC $\delta = 0.19$, both \ph{}) and the same density
($\delta = -0.09$, \regX{}). % src: STESTS outcomes files (full-aac-minimal-prose delta 0.121094) and loc (delta 0.19); TESTS secondary:violations_per_kloc (delta -0.085938)
On the registered primary field, \lev{full-aac} still scores higher by $+0.033$ ($\delta = 0.52$), and after the
LOC adjustment by $+0.034$ ($\delta = 0.47$, \ph{}). % src: TESTS primary:ahsCombined spec-level:full-aac-minimal-prose (0.032750, delta 0.519531); STESTS ahsCombined_resid_lnloc (SIZE §3: +0.034, delta 0.47)
The pairwise difference itself is a registered descriptive contrast; reading it as \emph{the} clean comparison is a
post-hoc choice made after the size result. It is the one specification contrast that size cannot explain, and it
favours the AaC specification. The vague prose specification is the worst level on the registered outcome.

\paragraph{Answer to RQ1.} Specification level changes the registered conformance score \regC{}. Against no
specification, both specification levels score lower, because they make the agent build more structure for the
rules to check \ph{}. Between the two specification styles at equal size, full AaC is better than minimal prose.

\subsection{RQ2: Model capability and vendor}\label{sec:rq2}

\paragraph{No model effect \regC{}.} The model main effect is not significant on $\AHSc$ (Holm $p = 0.438$) nor on
the co-primary $\AHSd$ (Holm $p = 0.753$). % src: TESTS families primary:ahsCombined (model p_holm 0.437956) and co-primary:ahsDeterministic (model p_holm 0.753225)
Mean $\AHSc$ differs by at most 0.014 between models (Haiku 0.910, Sonnet 0.920, Opus 0.924).
% src: TESTS primary:ahsCombined model pairs (largest |diff| 0.014030); SIZE §3 (model "all" rows)
Model size matters for whether a project type-checks (yield, above), not for the conformance of projects that do.
The interaction (Holm $p = 0.026$) reflects Haiku's different order (\lev{none} 0.929 $>$ \lev{minimal-prose} 0.910
$>$ \lev{full-aac} 0.895) against Opus and Sonnet (\lev{none} $>$ \lev{full-aac} $>$ \lev{minimal-prose}), on 3 or 4
Haiku cells per level; we do not interpret it further. % src: SIZE §3 (haiku-4-5 rows; opus, sonnet rows); GRID (Haiku valid cells 3/4/4)

\paragraph{The model effect lives only in the judge field \regX{}.} $\AHSn$ shows a large model effect (Holm
$p = 0.0003$; Haiku $-0.22$ against each of the others), while $\AHSd$ shows none. Under the registered rule, a model
effect seen only in the judge-bearing field is not read as a model effect. % src: TESTS family secondary:ahsNeuronal (model p_holm 0.000300; haiku-opus -0.222727); PLAN §6 (Co-primary)

\paragraph{The Codex arm reverses the order (exploratory).} On \texttt{gpt-5.6-terra}, \lev{minimal-prose} is best
(0.957) and \lev{none} worst (0.926) (Table~\ref{tab:levels}); the level test, \ph{} as a family of one, gives
$p = 0.028$. % src: SIZE §7; STESTS arm codex outcome ahsCombined spec-level (p_holm 0.027697)
Codex projects are smaller (6 to 13 files) and grow less with the specification. On \lev{none}, Codex loses AHS
through two SOLID proxies (\ff{FF-SO01} and \ff{FF-SO02}, each firing in 2 of 6 cells) and through judge integrity.
% src: SIZE §7
Per file, \lev{none} is the \emph{worst} Codex level (1.31 findings per file; \ph{} $p = 0.002$).
% src: SIZE §7 table; STESTS arm codex outcome violations_deterministic_per_file (p 0.0019)
The registered exploratory vendor contrast \regX{} favours \texttt{gpt-5.6-terra} over every Claude model on $\AHSc$
($p = 0.012$; differences $+0.018$ to $+0.032$, $\delta$ 0.22 to 0.61) and on $\AHSd$ ($p = 0.004$).
% src: TESTS families vendor:ahsCombined (p 0.012199; model:gpt-5.6-terra-* 0.017889..0.031919, delta 0.216..0.606) and vendor:ahsDeterministic (p 0.004100)
With one OpenAI model, vendor, model, harness, confinement and time are the same contrast, so this is not a vendor
effect in general. % src: PLAN §10 B10
The arm shows that the Claude-arm direction is not universal; with 6 cells per level it supports no claim.

\paragraph{Answer to RQ2.} Within the Claude family, model capability has no registered effect on conformance. The
exploratory second-vendor arm orders the levels differently and prefers minimal prose.

\subsection{RQ3: Failure families}\label{sec:rq3}

\paragraph{The registered profile is dominated by one rule \regX{}.} Counted by registered rule family on the Claude
arm, the symbolic findings are 1\,003 test-file-pairing findings (\ff{FF-CV05}, family FPAT-PATTERN), 154 coupling
findings and 51 SOLID-family findings: \ff{FF-CV05} makes 83\% of the 1\,208 findings.
% src: results/e1-grid/agg/so5_patterns.csv (basis = symbolic; derived sums by code over Claude run ids: 1003, 154, 51);
%      SIZE §5 ("FF-CV05 produces 1 003 of the 1 208 symbolic findings (83 %)")
It fires once per source file without a test file (0.96 findings per file on the Claude arm, 1.00 on Codex), so
raw counts mostly count files. % src: SIZE §5
The label-dependent family counts are not available, because the labelling population for E1 findings was dropped
before the run; the profile is therefore unlabelled and unweighted. % src: PLAN §6 (Family counts; "P3 is out of live labelling")

\paragraph{\ff{FF-CV05} does not move the AHS \ph{}.} \ff{FF-CV05} fires in all 47 Claude cells and costs each the
same 0.0125 of AHS, so it contributes nothing to any level contrast. % src: SFUNC rows FF-CV05 (cells_fired = n_valid; ahs_deficit_mean 0.012500 at every level)
This is the score property of Section~\ref{sec:instrument} at work: the AHS counts firing functions, not findings.

\paragraph{Two functions carry the gap \ph{}.} Splitting the AHS deficit by function
(Table~\ref{tab:decomp}, Figure~\ref{fig:firing}) shows that \ff{FF-C02} (module fan-out) and \ff{FF-P02}
(dependency inversion) carry most of every contrast, followed by \ff{FF-C06} (abstraction ratio).
\ff{FF-C02} fires in 0 of 15 \lev{none} cells, 15 of 16 \lev{minimal-prose} cells and 5 of 16 \lev{full-aac}
cells; \ff{FF-P02} fires in 2, 11 and 6. % src: SFUNC (cells_fired for FF-C02, FF-P02 by spec_level)
The judge dimensions show no level effect. % src: TESTS families secondary:fpat_FPAT-SEMANTIC (spec p_holm 1.0), secondary:fpat_FPAT-INTEGRITY (0.437), secondary:judge_fail_share (0.550)
The registered exploratory dimension tests agree: the level effect is significant on the coupling AVR (Holm
$p = 0.0003$) and the pattern AVR (Holm $p = 0.009$) \regX{}. % src: TESTS secondary:avr_coupling (spec p_holm 0.000300), secondary:avr_pattern (0.009199)
Within a level, firing does not track size: in \lev{full-aac}, cells where \ff{FF-C02} fired average 24 files and
cells where it did not average 32. % src: SIZE §5 (last paragraph)

\begin{table}[t]
\centering\small
\caption{Split of the $\AHSc$ gap by function, Claude arm \ph{}. Each entry is the function's contribution to the
mean AHS difference (``prose'' = \lev{minimal-prose}, ``AaC'' = \lev{full-aac}); the rows sum to the gap up to rounding.}
\label{tab:decomp}
\begin{tabular}{@{}llrrr@{}}
\toprule
Function & Dimension & prose $-$ none & AaC $-$ none & prose $-$ AaC \\
 & & (gap 0.064) & (gap 0.031) & (gap 0.033) \\
\midrule
\ff{FF-C02} module fan-out          & coupling   & $+0.028$ & $+0.009$ & $+0.019$ \\
\ff{FF-P02} dependency inversion    & pattern    & $+0.025$ & $+0.011$ & $+0.014$ \\
\ff{FF-C06} abstraction ratio       & coupling   & $+0.008$ & $+0.004$ & $+0.004$ \\
\ff{FF-C03} component instability   & coupling   & $+0.004$ & $+0.004$ & $0$ \\
\ff{FF-C04} orphan files            & coupling   & $0$      & $+0.002$ & $-0.002$ \\
\ff{FF-CV05} test-file pairing      & convention & $0$      & $0$      & $0$ \\
Judge integrity                     & integrity  & $0$      & $+0.002$ & $-0.002$ \\
\bottomrule
\end{tabular}
% src: SIZE §5 ("Contribution to the AHS gap" table), computed from SFUNC ahs_deficit_mean and
%      results/e1-grid/agg-size/size_dimensions.csv
\end{table}

\begin{figure}[t]
\centering
\begin{minipage}[t]{0.49\linewidth}\centering
% Generated by papers/paper1-e1/figures/make_figures.py from results/e1-grid/agg/so5_grid.csv. Do not edit by hand.
\definecolor{oiBlue}{HTML}{0072B2}
\definecolor{oiOrange}{HTML}{E69F00}
\definecolor{oiGreen}{HTML}{009E73}
\definecolor{oiPink}{HTML}{CC79A7}
\begin{tikzpicture}
\begin{axis}[width=\linewidth, height=5.0cm, xlabel={\texttt{.ts} files under \texttt{src/}}, ylabel={$\mathrm{AHS}_{\mathrm{comb}}$},
  ymin=0.84, ymax=1.0, xmin=0, ytick={0.85,0.90,0.95,1.00},
  yticklabel style={/pgf/number format/fixed,/pgf/number format/precision=2},
  grid=major, grid style={gray!20}, axis line style={gray!60},
  legend style={font=\scriptsize, at={(0.5,1.02)}, anchor=south, legend columns=3, draw=none},
  every axis plot/.append style={only marks}]
\addplot[mark=*, mark size=1.8pt, draw=oiBlue, fill=oiBlue, fill opacity=0.6] coordinates {(13,0.914) (8,0.884) (4,0.988) (10,0.957) (9,0.957) (9,0.957) (11,0.957) (15,0.957) (12,0.957) (8,0.957) (8,0.957) (8,0.957) (8,0.957) (9,0.957) (5,0.957)};
\addlegendentry{none}
\addplot[mark=square*, mark size=1.8pt, draw=oiOrange, fill=oiOrange, fill opacity=0.6] coordinates {(21,0.900) (16,0.912) (21,0.928) (22,0.900) (31,0.882) (40,0.882) (31,0.852) (34,0.882) (37,0.882) (32,0.882) (23,0.852) (24,0.852) (24,0.852) (27,0.928) (27,0.928) (29,0.882)};
\addlegendentry{minimal-prose}
\addplot[mark=triangle*, mark size=1.8pt, draw=oiGreen, fill=oiGreen, fill opacity=0.6] coordinates {(16,0.882) (15,0.929) (24,0.900) (23,0.869) (27,0.882) (46,0.957) (35,0.912) (41,0.957) (50,0.957) (37,0.957) (25,0.957) (25,0.912) (26,0.882) (28,0.957) (28,0.882) (28,0.928)};
\addlegendentry{full-aac}
\end{axis}
\end{tikzpicture}

\end{minipage}\hfill
\begin{minipage}[t]{0.49\linewidth}\centering
% Generated by papers/paper1-e1/figures/make_figures.py from results/e1-grid/agg-size/size_functions.csv. Do not edit by hand.
\definecolor{oiBlue}{HTML}{0072B2}
\definecolor{oiOrange}{HTML}{E69F00}
\definecolor{oiGreen}{HTML}{009E73}
\definecolor{oiPink}{HTML}{CC79A7}
\begin{tikzpicture}
\begin{axis}[width=\linewidth, height=5.0cm, ybar, bar width=4pt, ymin=0, ymax=1.08,
  symbolic x coords={FF-C02,FF-P02,FF-C06,FF-C03,FF-CV05}, xtick=data,
  xticklabels={{C02 fan-out},{P02 dep. inversion},{C06 abstraction},{C03 instability},{CV05 test pairing}}, xticklabel style={font=\tiny, align=center, text width=1.1cm},
  ylabel={share of cells that fire}, ytick={0,0.25,0.5,0.75,1},
  grid=major, ymajorgrids=true, xmajorgrids=false, grid style={gray!20}, axis line style={gray!60},
  legend style={font=\scriptsize, at={(0.5,1.02)}, anchor=south, legend columns=3, draw=none}]
\addplot[fill=oiBlue, draw=none, postaction={pattern=north east lines, pattern color=white}] coordinates {(FF-C02,0.0000) (FF-P02,0.1333) (FF-C06,0.0000) (FF-C03,0.8667) (FF-CV05,1.0000)};
\addlegendentry{none}
\addplot[fill=oiOrange, draw=none] coordinates {(FF-C02,0.9375) (FF-P02,0.6875) (FF-C06,0.2500) (FF-C03,1.0000) (FF-CV05,1.0000)};
\addlegendentry{minimal-prose}
\addplot[fill=oiGreen, draw=none, postaction={pattern=crosshatch, pattern color=white}] coordinates {(FF-C02,0.3125) (FF-P02,0.3750) (FF-C06,0.1250) (FF-C03,1.0000) (FF-CV05,1.0000)};
\addlegendentry{full-aac}
\end{axis}
\end{tikzpicture}

\end{minipage}
\caption{Left: file count against $\AHSc$ per valid Claude cell, by specification level (descriptive). Right: share
of valid Claude cells in which each function fires, by level \ph{}. Text alternative: \lev{none} cells have 4 to
15 files and high scores, \lev{minimal-prose} cells 16 to 40 and \lev{full-aac} cells 15 to 50; within
\lev{full-aac} larger projects score higher. \ff{FF-C02} fires in 0\%, 94\% and 31\% of cells at \lev{none},
\lev{minimal-prose} and \lev{full-aac}; \ff{FF-P02} in 13\%, 69\% and 38\%; \ff{FF-CV05} fires everywhere.
Generated by \texttt{figures/make\_figures.py} from \texttt{so5\_grid.csv} and \texttt{size\_functions.csv}.}
% src: file ranges derived from GRID (accepted Claude rows: none 4-15, minimal-prose 16-40, full-aac 15-50);
%      firing shares from SFUNC fire_rate (FF-C02 0, 0.9375, 0.3125; FF-P02 0.1333, 0.6875, 0.375)
\label{fig:files}\label{fig:firing}
\end{figure}

\paragraph{Answer to RQ3.} Module fan-out and dependency inversion drive the level differences. They need modules
and abstractions to fire; the small \lev{none} projects rarely give them any, and the AaC specification avoids them
more often than prose does. The most frequent finding, missing test files, is constant across levels and inert in the
score.

\subsection{RQ4: The LLM judge}\label{sec:rq4}

\paragraph{What the judge sees.} Across the 65 evaluated cells, the judge decided 300 integrity units (36 failing)
and 933 intent-alignment units (19 failing), with no invalid, split or unstable unit.
% src: JSUM rows FF-N01,* (300 valid, 36 fail) and FF-N02,* (933 valid, 19 fail; units_invalid, units_split_vote, units_unstable 0)
Every failing unit is on a file that a symbolic rule also flags. % src: JSUM column units_fail_rules_also_flag (36 and 19)
Fails concentrate on two generators: Haiku (33 of 209 valid units) and \texttt{gpt-5.6-terra} (18 of 230), against
Opus (2 of 450) and Sonnet (2 of 344). % src: derived, sums of units_fail and units_valid in JSUM by generator prefix of project_id
Under the registered coding frame, none of the seven judge criteria reaches the pre-declared discrimination
threshold of 0.3 when pooled (highest 0.258). % src: results/e1-grid/judge-diagnostics/judge_discrimination.csv (project_id *, max discrimination 0.258378, JC-N02-TRANSPORT)

\paragraph{At its registered weight the judge is inert \regX{}.} Under the registered weights ($W_n = 0.08$) the
judge changes no verdict and no ``judge mattered'' flag in any of the 65 cells, under the registered, proportional and
any-fail aggregations; its effect on $\AHSc$ ranges from $-0.029$ to 0 under the registered rule.
% src: AGGR (judge_mattered_registered, judge_mattered_proportional, verdict_changed, any_fail_changed: 0 of 65;
%      delta_judge_registered min -0.029, max 0); all rows reading = pre-registered
Only the most aggressive registered variant, judge weight 2/7 with any-fail aggregation, moves 8 cells from pass to
warning: 6 Haiku cells and 2 Codex cells. % src: AGGR judge_mattered_jw_any_fail (8 true: 6 claude-haiku-4-5/order-fulfilment, 2 gpt-5.6-terra)
The judge cannot explain the level effect of RQ1: the judge families show no level effect (Section~\ref{sec:rq3}).

\paragraph{Self-preference \regC{}.} The registered directional check is positive: mean $d$ is larger for
Opus-written cells than for the other Claude cells by $+0.083$ (one-sided $p = 0.014$, $\delta = 0.20$).
% src: TESTS effect directional:claude-opus-5-5-vs-others (statistic 0.083048, p 0.013699, cliffs_delta 0.197318)
The registered wording reports this as evidence of judge self-preference. The pattern behind it calls for care.
Mean $d$ is $+0.083$ for Opus, $+0.088$ for Sonnet and $-0.144$ for Haiku, and $\AHSn$ is 1.000 in every Opus and
Sonnet cell. % src: derived from GRID (mean of ahs_neuronal - ahs_deterministic over accepted rows by model: 0.083, 0.088, -0.144;
%      ahs_neuronal values: opus {1.000}, sonnet {1.000}, haiku {0.650, 1.000})
The contrast is therefore carried by the Haiku cells: the judge treats Opus and Sonnet code alike and finds more
problems in Haiku code than the rules do. The data cannot separate a preference for the judge's own model from a
judge that responds to lower-quality code, and Haiku's lower yield points to the second. The registered vendor check
is not significant ($+0.033$, $p = 0.173$, $\delta = 0.28$); by registration, that null is not evidence of no
cross-vendor self-preference. % src: TESTS effect directional:anthropic-vs-other-vendor (0.033035, p 0.172683, delta 0.275414); PLAN §6.1

\paragraph{Judge versus panel \regC{} / \regX{}.} The panel labelled all 46 sampled items (38 E1 units and 8 fixture
units) twice, with 92 calls, no retry and no tool use; replay is byte-identical. % src: aidlc-docs/audit.md "P4 labelled on the Gemini API" entry
Panel run-versus-run agreement is 0.957 [0.855, 0.988] ($\kappa = 0.777$, AC1 $= 0.955$), above the registered
$\kappa \geq 0.60$ taxonomy rule; the two runs disagree on 2 items, which become \texttt{uncertain}.
% src: P4 agreement row comparison run-vs-run (percentAgreement 0.9565, ci [0.8547, 0.9880], cohensKappa 0.7767, gwetAc1 0.9548, uncertain 2)
On the registered headline, the 36 E1 items with a decided panel label, weighted agreement is 0.982 [0.949, 1.000]
and AC1 is 0.982 [0.947, 1.000] (Table~\ref{tab:p4}). % src: P4 row judge-vs-panel source e1 headline true (n 36, 0.982, ci [0.949,1.000], AC1 0.982 [0.947,1.000])
Cohen's $\kappa$ is 0, because the judge said \emph{pass} on all 36 decided items: with one category on one side, $\kappa$ is
undefined in spirit and 0 by formula, the high-agreement paradox \citeneeded{feinstein1990paradox}; AC1 is the
registered companion for this case. % src: P4 (cohensKappa 0 on the headline row); aidlc-docs/audit.md ("The judge said pass on all 36 E1 pairs, so kappa is 0 by prevalence")
Both disagreements are the judge saying \emph{pass} where the panel said \emph{fail} on both runs, on
\ff{FF-N01} integrity units of \texttt{order-fulfilment}/\lev{none} projects (one Sonnet, one Codex). The judge's
errors in this sample are lenient. % src: aidlc-docs/audit.md (P4 entry: "The 2 disagreements are judge pass against panel fail (both runs) on FF-N01 order-fulfilment/none (sonnet, gpt)");
%      derived from P4 labels (sonnet item fail/fail, gpt item fail/fail)
On the 8 fixture items, which include 2 judge fails, judge and panel agree on all 8.
% src: P4 row source fixture (n 8, percentAgreement 1); aidlc-docs/audit.md ("Fixtures: n 8, all agree (2 fail, 6 pass)")

\begin{table}[t]
\centering\small
\setlength{\tabcolsep}{4pt}
\caption{Judge (\texttt{claude-opus-5-5}) versus panel (\texttt{gemini-3.1-pro-preview}). Weighted agreement with
weighted item-bootstrap intervals where the registered rule computes one; ``counts'' rows have too few items.}
\label{tab:p4}
\begin{tabular}{@{}lrrrr@{}}
\toprule
Row & $n$ & Agreement & $\kappa$ & AC1 \\
\midrule
E1, headline \regC{}                      & 36 & 0.982 [0.949, 1.000] & 0.000 & 0.982 [0.947, 1.000] \\
E1, \texttt{uncertain} as a category \regX{} & 38 & 0.965 [0.920, 0.993] & 0.192 & 0.964 [0.916, 0.993] \\
Fixtures \regX{}                           & 8  & 1.000 (counts) & 1.000 & 1.000 \\
Pooled \regX{}                             & 44 & 0.984 [0.956, 1.000] & 0.742 & 0.983 [0.951, 1.000] \\
\midrule
E1, Haiku-written \regX{}                  & 8  & 1.000 (counts) & n/a & 1.000 \\
E1, Opus-written \regX{}                   & 9  & 1.000 (counts) & n/a & 1.000 \\
E1, Sonnet-written \regX{}                 & 9  & 0.968 (counts) & 0.000 & 0.967 \\
E1, GPT-written \regX{}                    & 10 & 0.953 [0.809, 1.000] & 0.000 & 0.951 [0.770, 1.000] \\
Panel run vs run \regX{}                   & 46 & 0.957 [0.855, 0.988] & 0.777 & 0.955 [0.886, 1.000] \\
\bottomrule
\end{tabular}
% src: P4 agreement rows (judge-vs-panel: source e1 headline; source e1 uncertainAsCategory true; source fixture;
%      source all; per generatorModel; run-vs-run). Per-model rows use ciMethod counts-only except gpt-5.6-terra.
\end{table}

\paragraph{What the agreement does not show.} The P4 frame holds 425 E1 judge verdicts, 20 of them \emph{fail}.
The registered stratified draw took 38 E1 units, of which the judge failed one, and on that unit the two panel runs
split, so it is \texttt{uncertain} and leaves the headline. % src: derived from results/e1-grid/label-plan-p4/judge-verdicts.json
% (source e1: 405 pass, 20 fail) joined to P4 labels on (projectId, functionId, unitId): E1 judge pass 36 decided + 1 uncertain, judge fail 1 uncertain
The headline therefore validates the judge's passes on E1 code and says nothing about its fails there; on fixtures its
2 fails were confirmed. The planned half-width for the E1 headline was $\pm 0.186$ at 50\% agreement, so the interval
is narrow only because agreement is near 1. % src: results/e1-grid/label-plan-p4/label-plan-summary.json (precision row "P4 judge-vs-panel agreement, E1 headline": halfWidthAt50 0.186)

\paragraph{Robustness to the judge model \ph{}.} A separate check, registered after the E1 results were known but
before any of its calls (ADR-038), replayed 100 recorded judge units, 50 of them from E1 cells, through the same
pinned route with only the model changed from \texttt{claude-opus-5-5} to \texttt{claude-sonnet-5-5}; each unit was
again judged three times. Over all 100 units the two models agree on 94\% of unit verdicts, Cohen's
$\kappa = 0.766$ [0.565, 0.926] and AC1 0.931, above the pre-declared floor of 0.60 for reading the verdicts as robust
to the judge model. % src: ADR.md ADR-038 Decision items 1-4; results/judge-robustness/agreement.csv (row all, primary: agreement 0.940000, kappa 0.766355 [0.565170, 0.925651], ac1 0.931232); results/judge-robustness/summary.json (units 100, calls 300, interpretation "model-robust (kappa >= 0.60)")
On the 50 E1 units, which include 6 Opus fails, the two models agree on every verdict; all six disagreements are on
real-project units where Opus failed and Sonnet passed, so Opus is slightly the stricter judge (18 against 12 failing
units of 100). Substituting the Sonnet votes changes no function, project or cell verdict in any of the 41 affected
reports. % src: results/judge-robustness/agreement.csv (E1 row: n 50, agreement 1.000000); results/judge-robustness/units.csv (E1: 6 fail/fail, 44 pass/pass, derived); results/judge-robustness/transition.csv (opus fail / sonnet pass 6); results/judge-robustness/fail-rates.csv (18 and 12); results/judge-robustness/summary.json (reportsAffected 41, functionVerdictChanges 0, registeredVerdictChanges 0, anyFailVerdictChanges 0)
This is agreement between two models of one vendor, a reliability check, not validation: it shows that the judge's
null effect on E1 verdicts does not hinge on the model chosen, not that the judge is right.

\paragraph{Answer to RQ4.} At an 8\% weight the judge adds diagnostic text but no verdict change, and it does not
contribute to the specification effect. The registered self-preference check is positive, but it is carried by the
weakest model's cells and confounded with code quality. A panel from another vendor agrees with the judge on 34 of 36
sampled E1 units (weighted agreement 0.982), and where they disagree the judge is the lenient one. Replacing the judge model changes no verdict.

% =====================================================================================================================
\section{Discussion}\label{sec:discussion}

\subsection{What specifications do to LLM architecture}

The registered result, taken alone, says that a specification makes generated architecture worse. We think that
reading is wrong, and the data show why. A specification changes \emph{how much} architecture the agent builds.
Without one, the agents produce a working service in about nine files: about three per layer, few interfaces,
few modules that import many others. With a specification, they produce three times as many files, with repository
interfaces, use cases, DTOs and adapters, as the specification asks. This richer structure gives the coupling and
dependency-inversion rules places to fire; a nine-file project gives them almost nothing to check. Per file, the
generated code is equally conformant at every level. The conformance score registers the change in structure, and
because it rewards structures that trigger fewer rules, it ranks the smallest projects first. That a small project
scores well does not show that its design is better: the \lev{none} projects contain the fewest of the elements that
clean architecture prescribes (repository interfaces, use cases, adapters), and so the fewest places where its rules
can be checked.

The comparison that size cannot explain is between the two specification styles. They produce projects of the same
size and the same finding density, yet the full AaC specification leads to fewer firings of module fan-out and
dependency inversion. The likely mechanism is that the AaC document names these properties as checkable functions
with thresholds, while the prose names the layers, the dependency direction and coding against abstractions,
but says nothing about fan-out or abstraction ratios and gives no number. % src: scripts/generator/prompts/minimal-prose.md (architecture paragraph) That explanation is a hypothesis: at \lev{full-aac} the agent sees the evaluator's own rules, so the
advantage could be ``knowing the test'' rather than ``understanding the architecture'' (Section~\ref{sec:threats}).
Either way, for a team that evaluates generated code against an AaC specification, giving the agent that
specification helps, and giving it a vague paraphrase is the worst option we tested.

The Codex arm is a warning against generalising the direction. Its agent builds smaller projects that grow less with
the specification, scores worst without one and best with prose. With one model, 18 cells and a different harness,
we cannot say whether this is a property of the vendor, of the model or of the harness.

Model capability did not matter for conformance among projects that compile. It did matter for whether they
compile: Haiku lost 7 of 18 cells to the type-check. A practitioner choosing between Claude tiers for this kind of
task gains reliability, not architecture, from the larger models.

\subsection{Measuring generated architecture}

E1 has two lessons for anyone who scores generated code. First, a score built from rule firings is not
size-neutral: when a treatment changes size, report size beside the score and add a per-file reading. Second, a
rule that fires everywhere (here, missing test files) dominates finding counts while contributing nothing to
contrasts; counts by family mostly count files. Both were visible only because the analysis kept per-function
detail, and both were found after the registered tests. We report them as post-hoc and recommend that future
designs register a size-normalised and a per-file outcome in advance.

On the judge, E1 replicates the companion study: with an 8\% weight and a majority rule, the judge cannot move a
verdict in a pass band this wide \cite{daedalus2026tool}. A judge earns its place in such a score only if its weight
is large enough to matter and its fails are validated; E1 validates its passes but not its fails. Its verdicts do
not hinge on the judge model: a second Claude model agrees on every sampled E1 unit and changes no verdict.
% src: results/judge-robustness/agreement.csv (E1 row: n 50, agreement 1.000000); results/judge-robustness/summary.json (registeredVerdictChanges 0)

\subsection{Practical guidance: what to put in a specification}

These recommendations follow from E1's evidence and are bounded by its scope (two Express tasks, mostly one vendor,
clean architecture). They are hypotheses for practice, not established rules.
\begin{enumerate}
  \item \textbf{If you will check the architecture, give the agent the checkable specification.} The full AaC
        document gave more conformant code than prose at equal size (\regC{} contrast, post-hoc reading). A paraphrase
        that names layers but not the measured properties was the worst level.
  \item \textbf{Name the coupling and inversion properties you care about, with limits.} The gap between prose and
        AaC sits in module fan-out and dependency inversion. The prose asked agents to code against abstractions
        but gave no fan-out limit and no check. A prose specification
        should at least say that use cases depend on interfaces, that no module imports many others, and what the
        limit is.
  \item \textbf{Expect a specification to make the project bigger.} A layered specification roughly triples the
        file count here. Budget review effort and evaluation accordingly, and do not compare scores across projects
        of different size without a per-file view.
  \item \textbf{Ask for tests explicitly if you want them.} The test-file-pairing rule fired in every cell, about
        once per source file, even at \lev{full-aac}, whose preset lists that rule. The task briefs and the prose
        level never mention tests; listing a rule among 27 was not enough to make the agents write them.
  \item \textbf{Do not expect a larger model to fix architecture.} Within one vendor, a larger model raised the
        chance of a compiling project, not its conformance.
\end{enumerate}

% =====================================================================================================================
\section{Threats to Validity}\label{sec:threats}

We follow the registered threats register (\texttt{Docs/threats-to-validity.md}) and the analysis plan's reporting
duties; the most important threats are below \citeneeded{wohlin2012experimentation,sallou2024breaking}.

\paragraph{Construct validity.}
\emph{The \lev{full-aac} level shows the evaluator specification.} The agent receives a preset with the same rules,
thresholds and weights the instrument uses, so the level is confounded with knowledge of the test and with prompt
length (21 to 637 template lines). % src: PLAN §10 B5; GENP §5; derived wc -l of the templates
The size-matched AaC advantage may partly measure ``teaching to the test''. A fourth level that gives an equally
long, equally precise specification of a \emph{different} rule set would separate the two.
\emph{The AHS rewards small structures} (Section~\ref{sec:discussion}); we report per-file and size-adjusted readings,
all post-hoc. \emph{Instrument precision on real code is uncertain}: 0.328 to 0.789 depending on how
\texttt{uncertain} labels count, with wide intervals; for the test-file-pairing rule, 2 of 6 sampled real findings
were judged real and 3 were uncertain, while all 3 sampled dependency-inversion findings were judged real.
% src: results/so4-heldout/agg/precision_baseline.csv (v2 overall 0.327901 [0.132388, 0.609359]; function FF-CV05 n 6, tp 2, uncertain 3;
%      function FF-P02 n 3, tp 3); ADR.md ADR-033 item 4 (0.608, 0.789)
Those figures come from real projects, not from E1 code, and E1 findings were not labelled.
\emph{Circularity.} The specification, the rules, the prompt levels and the instrument were written by the same
author, and the prose level was derived from the same preset. The study measures conformance to this author's
rendering of clean architecture. % src: TTV §2.8; Docs/reports/director-report-2026-10.md §5 item 1
\emph{All cells pass}: the study measures differences inside the pass band, never a blocking violation.

\paragraph{Internal validity.}
\emph{The ADR-032 fix.} The first judge pass found that the extractor had no registered way to read an E1 tree's
tsconfig; 9 Haiku trees held an agent-written one and 63 held none. The fix, decided during the run, makes the
instrument read the pinned per-run tsconfig of the generator protocol for every cell and ignore any agent-written one;
all pass-1 records were discarded and the evaluation re-ran on every cell under prereg v20. It changes how files are
located, not any rule, weight or threshold, but it is a post-hoc operational deviation.
% src: ADR.md ADR-032 (Context, Decision items 1-4, Classification); GENP §3 (Evaluation note)
\emph{Usage-limit handling.} In the re-run, one usage window was classified as a transport error instead of the
registered pause; the same command was re-run with cassettes reused and no cell was lost.
% src: aidlc-docs/audit.md ("go E1" entry); TTV TV-38
\emph{Seven Haiku cells were not run} because their generations failed the type-check. The failed cells are not
random: the surviving Haiku cells are the ones that compiled, so Haiku's conformance is measured on a selected
subset (3 or 4 cells per level). % src: GRID; SIZE §3 (haiku valid/coords 3/6, 4/6, 4/6)
\emph{Judge self-preference} is confounded with code quality (Section~\ref{sec:rq4}).

\paragraph{Statistical conclusion validity.}
\emph{$n = 3$ per cell.} Three replicates of a stochastic generator, no temperature or seed control, and 47 valid
Claude cells in all. The pairwise intervals resample three runs within each cell and are descriptive only.
% src: PLAN §6 (Effect sizes); TTV TV-78
\emph{Post-hoc families.} The size analysis adds many tests, each in its own Holm family, with no correction across
families, chosen after the registered results were known. % src: SIZE §3, §8
\emph{The judge-versus-panel headline} rests on 36 items with no judge fail among them (the one sampled judge fail
was \texttt{uncertain}); $\kappa$ is uninformative and
the judge's fails on E1 code are not validated. The model-robustness check compares two Claude models and is
reliability evidence only.

\paragraph{External validity.}
\emph{Two tasks}, both small in-memory Express services, as a blocking factor, not a sample of tasks.
\emph{Single-vendor confirmatory tests}: every registered test is on Claude models judged by a Claude model; the
Codex arm is one model, exploratory, and differs in harness, confinement, reasoning settings and time.
% src: PLAN §10 B5, B10; TTV TV-80, TV-100
\emph{One architecture style} (clean architecture) and one language. \emph{Model generation} is confounded with tier
(Haiku 4.5 against two 5.5 models). % src: PLAN §10 B5
The models and CLIs are pinned versions that will change.

\paragraph{Reliability.} Every generated tree is identified by its git tree hash, every judge and labeller call is
stored as a cassette and replays without a model call, and the aggregation is byte-identical across reruns.
% src: GENP §7 (treeSha); TTV TV-81; SIZE header ("Two runs give the same bytes")
The size analysis reproduces every registered row it overlaps.

% =====================================================================================================================
\section{Related Work}\label{sec:related}

\paragraph{Quality of LLM-generated code.} Benchmarks such as HumanEval and SWE-bench measure functional correctness
of generated code \citeneeded{chen2021codex,jimenez2024swebench}, and stricter test suites show that correctness is
overestimated \citeneeded{liu2023evalplus}. Studies of generated code's maintainability and smells exist
\citeneeded{siddiq2024quality}; we are not aware of a registered factorial study of architectural conformance of
whole agent-generated services. \todo{search 2025--2026 work on architecture-level quality of agent-generated code
and on specification or prompt form as a factor; position against it.}

\paragraph{LLMs and software architecture.} Work on LLMs for architecture has studied design decision generation
\citeneeded{dhar2024adr} and the state of practice \citeneeded{jahic2024practice}. Spec-driven tools make the
specification an input to the agent \citeneeded{speckit2025,kiro2025}; prompt-pattern catalogues treat the prompt's
form as a design choice \citeneeded{white2023prompt}. E1 is, to our knowledge, the first test of whether the
specification's form changes architectural conformance under pre-registration. \todo{verify this priority claim.}

\paragraph{Architecture conformance checking.} Reflexion models and their successors compare intended and implemented
structure \citeneeded{murphy1995reflexion,koschke2003hierarchical}; surveys compare static approaches
\citeneeded{knodel2007comparison,passos2010static}, and erosion is a well-documented cost
\citeneeded{desilva2012erosion,rosik2011drift}. Fitness functions make conformance continuous
\citeneeded{ford2017evolutionary}. Our instrument and its validation are in the companion paper
\cite{daedalus2026tool}.

\paragraph{LLM-as-a-judge.} LLM judges are biased by position and verbosity \citeneeded{wang2023fair,zheng2023judging}
and can favour their own outputs \citeneeded{panickssery2024self,wataoka2024selfpref}. Our self-preference check
follows that concern; our result shows how such a check can be confounded with quality when the judge's family also
contains the weakest generator.

\paragraph{Method.} We follow guidance on experimentation and empirical standards in software engineering
\citeneeded{kitchenham2002guidelines,wohlin2012experimentation,ralph2020standards} and on pre-registration
\citeneeded{nosek2018prereg}; threats specific to LLM-based studies follow \citeneeded{sallou2024breaking}.

% =====================================================================================================================
\section{Conclusion and Open Data}\label{sec:conclusion}

In a pre-registered experiment with three Claude models, three specification levels and two tasks, the
specification level changed the architectural conformance of agent-generated TypeScript services, and the model did
not. The registered direction favoured no specification, but a post-hoc analysis traced this to project size: with a
specification, agents build about three times as many files, and the two rules that need structure, module fan-out
and dependency inversion, then fire. Findings per file did not differ. At equal size, a full AaC specification gave
better conformance than a short prose one, which was the worst level. An exploratory second-vendor arm reversed the
order of levels. The LLM judge, at its 8\% weight, changed no verdict; a positive self-preference check was carried
by the weakest model's cells; and a cross-vendor panel agreed with the judge's passes on 34 of 36 sampled units (weighted agreement 0.982).

Next steps are a registered replication with a size-normalised primary outcome, a fourth level that gives a precise
specification of a different rule set to separate understanding from knowing the test, more tasks and vendors, and a
judge-validation sample stratified on judge fails.

\paragraph{Open data.} The replication package contains the registered plans, generator protocol, prompt templates
and seeds (\texttt{experiments/e1-grid/}, \texttt{Docs/}), the hash ledger (\texttt{corpus/prereg.json}), every
evaluation record and aggregate table (\texttt{results/e1-grid/}), all judge and labeller cassettes, the post-hoc size
analysis and its producer (\texttt{scripts/e1-size-analysis-cli.ts}), and the script that generates this paper's
figures from the committed CSVs (\texttt{papers/paper1-e1/figures/make\_figures.py}). Every number in this paper
carries a source-path comment in the \LaTeX{} source. The generated project trees live outside the repository and are
identified by their tree hashes. \todo{archive DOI; decide whether to publish the generated trees.}

\bibliographystyle{plain}
\bibliography{refs}

\end{document}
